\documentclass[10pt,a4paper]{amsart}
\usepackage{amsmath,amssymb,amsthm,graphicx}
% For commutative diagrams we use Xy-pic
\usepackage[all]{xy}

% We use 2cell for 2-commutative diagrams.
\xyoption{2cell}
% Theorem environments.
%
\theoremstyle{plain}
\newtheorem{theorem}[subsection]{Theorem}
\newtheorem{proposition}[subsection]{Proposition}
\newtheorem{lemma}[subsection]{Lemma}

\theoremstyle{definition}
\newtheorem{definition}[subsection]{Definition}
\newtheorem{example}[subsection]{Example}
\newtheorem{exercise}[subsection]{Exercise}
\newtheorem{situation}[subsection]{Situation}

\theoremstyle{remark}
\newtheorem{remark}[subsection]{Remark}
\newtheorem{remarks}[subsection]{Remarks}

\numberwithin{equation}{subsection}

% Macros
%
\def\lim{\mathop{\mathrm{lim}}\nolimits}
\def\colim{\mathop{\mathrm{colim}}\nolimits}
\def\Spec{\mathop{\mathrm{Spec}}}
\def\Hom{\mathop{\mathrm{Hom}}\nolimits}
\def\Ext{\mathop{\mathrm{Ext}}\nolimits}
\def\SheafHom{\mathop{\mathcal{H}\!\mathit{om}}\nolimits}
\def\SheafExt{\mathop{\mathcal{E}\!\mathit{xt}}\nolimits}
\def\Sch{\mathit{Sch}}
\def\Mor{\mathop{\mathrm{Mor}}\nolimits}
\def\Ob{\mathop{\mathrm{Ob}}\nolimits}
\def\Sh{\mathop{\mathit{Sh}}\nolimits}
\def\NL{\mathop{N\!L}\nolimits}
\def\CH{\mathop{\mathrm{CH}}\nolimits}
\def\proetale{{pro\text{-}\acute{e}tale}}
\def\etale{{\acute{e}tale}}
\def\QCoh{\mathit{QCoh}}
\def\Ker{\mathop{\mathrm{Ker}}}
\def\Im{\mathop{\mathrm{Im}}}
\def\Coker{\mathop{\mathrm{Coker}}}
\def\Coim{\mathop{\mathrm{Coim}}}

% Boxtimes
%
\DeclareMathSymbol{\boxtimes}{\mathbin}{AMSa}{"02}

%
% Macros for moduli stacks/spaces
%
\def\QCohstack{\mathcal{QC}\!\mathit{oh}}
\def\Cohstack{\mathcal{C}\!\mathit{oh}}
\def\Spacesstack{\mathcal{S}\!\mathit{paces}}
\def\Quotfunctor{\mathrm{Quot}}
\def\Hilbfunctor{\mathrm{Hilb}}
\def\Curvesstack{\mathcal{C}\!\mathit{urves}}
\def\Polarizedstack{\mathcal{P}\!\mathit{olarized}}
\def\Complexesstack{\mathcal{C}\!\mathit{omplexes}}
% \Pic is the operator that assigns to X its picard group, usage \Pic(X)
% \Picardstack_{X/B} denotes the Picard stack of X over B
% \Picardfunctor_{X/B} denotes the Picard functor of X over B
\def\Pic{\mathop{\mathrm{Pic}}\nolimits}
\def\Picardstack{\mathcal{P}\!\mathit{ic}}
\def\Picardfunctor{\mathrm{Pic}}
\def\Deformationcategory{\mathcal{D}\!\mathit{ef}}
\newenvironment{reference}{}{}
\newenvironment{slogan}{}{}
\newenvironment{history}{}{}
\usepackage{cleveref}
\crefname{originallemma}{original lemma}{original lemmas}
\Crefname{originallemma}{Original lemma}{Original lemmas}
\makeatletter
\newlabel{quot-lemma-polarized-limits}{{0D43}{source}{}{sourcepointer-1}{}}
\newlabel{quot-lemma-polarized-limits@cref}{{[originallemma][0][]0D43}{[1][0][]source}{}{}{}}
\makeatother
\crefname{originallemma}{original lemma}{original lemmas}
\Crefname{originallemma}{Original lemma}{Original lemmas}
\makeatletter
\newlabel{stacks-limits-lemma-limit-preserving-diagonal}{{0CMW}{source}{}{sourcepointer-2}{}}
\newlabel{stacks-limits-lemma-limit-preserving-diagonal@cref}{{[originallemma][0][]0CMW}{[1][0][]source}{}{}{}}
\makeatother
\crefname{originalproposition}{original proposition}{original propositions}
\Crefname{originalproposition}{Original proposition}{Original propositions}
\makeatletter
\newlabel{stacks-limits-proposition-characterize-locally-finite-presentation}{{0CMY}{source}{}{sourcepointer-3}{}}
\newlabel{stacks-limits-proposition-characterize-locally-finite-presentation@cref}{{[originalproposition][0][]0CMY}{[1][0][]source}{}{}{}}
\makeatother
\crefname{originallemma}{original lemma}{original lemmas}
\Crefname{originallemma}{Original lemma}{Original lemmas}
\makeatletter
\newlabel{lemma-Mor-s-lfp}{{0DPN}{source}{}{sourcepointer-4}{}}
\newlabel{lemma-Mor-s-lfp@cref}{{[originallemma][0][]0DPN}{[1][0][]source}{}{}{}}
\makeatother
\crefname{originallemma}{original lemma}{original lemmas}
\Crefname{originallemma}{Original lemma}{Original lemmas}
\makeatletter
\newlabel{lemma-Isom-in-Mor}{{0DPP}{source}{}{sourcepointer-5}{}}
\newlabel{lemma-Isom-in-Mor@cref}{{[originallemma][0][]0DPP}{[1][0][]source}{}{}{}}
\makeatother
\crefname{originalsection}{original section}{original sections}
\Crefname{originalsection}{Original section}{Original sections}
\makeatletter
\newlabel{section-hom-isom}{{0DLW}{source}{}{sourcepointer-6}{}}
\newlabel{section-hom-isom@cref}{{[originalsection][0][]0DLW}{[1][0][]source}{}{}{}}
\makeatother
\crefname{originallemma}{original lemma}{original lemmas}
\Crefname{originallemma}{Original lemma}{Original lemmas}
\makeatletter
\newlabel{stacks-properties-lemma-check-property-covering}{{04XD}{source}{}{sourcepointer-7}{}}
\newlabel{stacks-properties-lemma-check-property-covering@cref}{{[originallemma][0][]04XD}{[1][0][]source}{}{}{}}
\makeatother
\crefname{originallemma}{original lemma}{original lemmas}
\Crefname{originallemma}{Original lemma}{Original lemmas}
\makeatletter
\newlabel{stacks-morphisms-lemma-composition-finite-presentation}{{06Q3}{source}{}{sourcepointer-8}{}}
\newlabel{stacks-morphisms-lemma-composition-finite-presentation@cref}{{[originallemma][0][]06Q3}{[1][0][]source}{}{}{}}
\makeatother
\crefname{originallemma}{original lemma}{original lemmas}
\Crefname{originallemma}{Original lemma}{Original lemmas}
\makeatletter
\newlabel{stacks-morphisms-lemma-smooth-locally-finite-presentation}{{0DNP}{source}{}{sourcepointer-9}{}}
\newlabel{stacks-morphisms-lemma-smooth-locally-finite-presentation@cref}{{[originallemma][0][]0DNP}{[1][0][]source}{}{}{}}
\makeatother
\crefname{originallemma}{original lemma}{original lemmas}
\Crefname{originallemma}{Original lemma}{Original lemmas}
\makeatletter
\newlabel{varieties-lemma-numerical-polynomial-from-euler}{{0BEM}{source}{}{sourcepointer-10}{}}
\newlabel{varieties-lemma-numerical-polynomial-from-euler@cref}{{[originallemma][0][]0BEM}{[1][0][]source}{}{}{}}
\makeatother
\crefname{originallemma}{original lemma}{original lemmas}
\Crefname{originallemma}{Original lemma}{Original lemmas}
\makeatletter
\newlabel{perfect-lemma-chi-locally-constant-geometric}{{0B9T}{source}{}{sourcepointer-11}{}}
\newlabel{perfect-lemma-chi-locally-constant-geometric@cref}{{[originallemma][0][]0B9T}{[1][0][]source}{}{}{}}
\makeatother
\crefname{originalsection}{original section}{original sections}
\Crefname{originalsection}{Original section}{Original sections}
\makeatletter
\newlabel{remark-Mor-numerical}{{0DPQ}{source}{}{sourcepointer-12}{}}
\newlabel{remark-Mor-numerical@cref}{{[originalsection][0][]0DPQ}{[1][0][]source}{}{}{}}
\makeatother
\crefname{originallemma}{original lemma}{original lemmas}
\Crefname{originallemma}{Original lemma}{Original lemmas}
\makeatletter
\newlabel{lemma-Mor-qc-over-base}{{0DPR}{source}{}{sourcepointer-13}{}}
\newlabel{lemma-Mor-qc-over-base@cref}{{[originallemma][0][]0DPR}{[1][0][]source}{}{}{}}
\makeatother
\crefname{originaltheorem}{original theorem}{original theorems}
\Crefname{originaltheorem}{Original theorem}{Original theorems}
\makeatletter
\newlabel{quot-theorem-polarized-algebraic}{{0D4X}{source}{}{sourcepointer-14}{}}
\newlabel{quot-theorem-polarized-algebraic@cref}{{[originaltheorem][0][]0D4X}{[1][0][]source}{}{}{}}
\makeatother
\crefname{originaldefinition}{original definition}{original definitions}
\Crefname{originaldefinition}{Original definition}{Original definitions}
\makeatletter
\newlabel{morphisms-definition-relatively-ample}{{01VH}{source}{}{sourcepointer-15}{}}
\newlabel{morphisms-definition-relatively-ample@cref}{{[originaldefinition][0][]01VH}{[1][0][]source}{}{}{}}
\makeatother
\crefname{originalsection}{original section}{original sections}
\Crefname{originalsection}{Original section}{Original sections}
\makeatletter
\newlabel{examples-stacks-section-stack-of-spaces}{{04SP}{source}{}{sourcepointer-16}{}}
\newlabel{examples-stacks-section-stack-of-spaces@cref}{{[originalsection][0][]04SP}{[1][0][]source}{}{}{}}
\makeatother
\crefname{originalsection}{original section}{original sections}
\Crefname{originalsection}{Original section}{Original sections}
\makeatletter
\newlabel{examples-stacks-section-stack-of-quasi-coherent-sheaves}{{03YL}{source}{}{sourcepointer-17}{}}
\newlabel{examples-stacks-section-stack-of-quasi-coherent-sheaves@cref}{{[originalsection][0][]03YL}{[1][0][]source}{}{}{}}
\makeatother
\crefname{originalsection}{original section}{original sections}
\Crefname{originalsection}{Original section}{Original sections}
\makeatletter
\newlabel{section-conventions}{{05X6}{source}{}{sourcepointer-18}{}}
\newlabel{section-conventions@cref}{{[originalsection][0][]05X6}{[1][0][]source}{}{}{}}
\makeatother
\title{Separated finite-presentation diagonal of the stack of polarized proper schemes}
\author{The Stacks Project Authors}
\date{stacks-project:moduli.tex:local-snapshot-20261011; GNU FDL 1.2 or later}
\begin{document}
\maketitle
\noindent\textit{Editorial note.} The native definition ofPolarized is included fromquot.tex3765–3800: proper flat finitely presented schemes with a relatively ample invertible sheaf and polarization-preserving isomorphisms.\par\medskip
\noindent\textit{Allowed standard prerequisites (editorial).}\begin{itemize}
\item Explicit prerequisite boundary. Proper flat finitely presented families, relatively ample line bundles, representability of Hom/Isom and polarized moduli, cohomology/base change for large tensor powers, projective embeddings and limit criteria for finite presentation.
\end{itemize}\medskip
\ifdefined\hypertarget\hypertarget{sourcepointer-1}{}\fi
\noindent\textit{Source reference.} Original lemma 0D43: focused-seeds/AGP/sources/stacks-snapshot-20261011/tags.txt, source lines 16438--16438. Stable Stacks tag for an expressly allowed foundational or notational input; the counted construction remains the two concrete core obligations. The retained local core is not replaced by this pointer. See standard\_prerequisites.\par\medskip
\ifdefined\hypertarget\hypertarget{sourcepointer-2}{}\fi
\noindent\textit{Source reference.} Original lemma 0CMW: focused-seeds/AGP/sources/stacks-snapshot-20261011/tags.txt, source lines 15871--15871. Stable Stacks tag for an expressly allowed foundational or notational input; the counted construction remains the two concrete core obligations. The retained local core is not replaced by this pointer. See standard\_prerequisites.\par\medskip
\ifdefined\hypertarget\hypertarget{sourcepointer-3}{}\fi
\noindent\textit{Source reference.} Original proposition 0CMY: focused-seeds/AGP/sources/stacks-snapshot-20261011/tags.txt, source lines 15873--15873. Stable Stacks tag for an expressly allowed foundational or notational input; the counted construction remains the two concrete core obligations. The retained local core is not replaced by this pointer. See standard\_prerequisites.\par\medskip
\ifdefined\hypertarget\hypertarget{sourcepointer-4}{}\fi
\noindent\textit{Source reference.} Original lemma 0DPN: focused-seeds/AGP/sources/stacks-snapshot-20261011/tags.txt, source lines 17158--17158. Stable Stacks tag for an expressly allowed foundational or notational input; the counted construction remains the two concrete core obligations. The retained local core is not replaced by this pointer. See standard\_prerequisites.\par\medskip
\ifdefined\hypertarget\hypertarget{sourcepointer-5}{}\fi
\noindent\textit{Source reference.} Original lemma 0DPP: focused-seeds/AGP/sources/stacks-snapshot-20261011/tags.txt, source lines 17159--17159. Stable Stacks tag for an expressly allowed foundational or notational input; the counted construction remains the two concrete core obligations. The retained local core is not replaced by this pointer. See standard\_prerequisites.\par\medskip
\ifdefined\hypertarget\hypertarget{sourcepointer-6}{}\fi
\noindent\textit{Source reference.} Original section 0DLW: focused-seeds/AGP/sources/stacks-snapshot-20261011/tags.txt, source lines 17061--17061. Stable Stacks tag for an expressly allowed foundational or notational input; the counted construction remains the two concrete core obligations. The retained local core is not replaced by this pointer. See standard\_prerequisites.\par\medskip
\ifdefined\hypertarget\hypertarget{sourcepointer-7}{}\fi
\noindent\textit{Source reference.} Original lemma 04XD: focused-seeds/AGP/sources/stacks-snapshot-20261011/tags.txt, source lines 6391--6391. Stable Stacks tag for an expressly allowed foundational or notational input; the counted construction remains the two concrete core obligations. The retained local core is not replaced by this pointer. See standard\_prerequisites.\par\medskip
\ifdefined\hypertarget\hypertarget{sourcepointer-8}{}\fi
\noindent\textit{Source reference.} Original lemma 06Q3: focused-seeds/AGP/sources/stacks-snapshot-20261011/tags.txt, source lines 8592--8592. Stable Stacks tag for an expressly allowed foundational or notational input; the counted construction remains the two concrete core obligations. The retained local core is not replaced by this pointer. See standard\_prerequisites.\par\medskip
\ifdefined\hypertarget\hypertarget{sourcepointer-9}{}\fi
\noindent\textit{Source reference.} Original lemma 0DNP: focused-seeds/AGP/sources/stacks-snapshot-20261011/tags.txt, source lines 17124--17124. Stable Stacks tag for an expressly allowed foundational or notational input; the counted construction remains the two concrete core obligations. The retained local core is not replaced by this pointer. See standard\_prerequisites.\par\medskip
\ifdefined\hypertarget\hypertarget{sourcepointer-10}{}\fi
\noindent\textit{Source reference.} Original lemma 0BEM: focused-seeds/AGP/sources/stacks-snapshot-20261011/tags.txt, source lines 14355--14355. Stable Stacks tag for an expressly allowed foundational or notational input; the counted construction remains the two concrete core obligations. The retained local core is not replaced by this pointer. See standard\_prerequisites.\par\medskip
\ifdefined\hypertarget\hypertarget{sourcepointer-11}{}\fi
\noindent\textit{Source reference.} Original lemma 0B9T: focused-seeds/AGP/sources/stacks-snapshot-20261011/tags.txt, source lines 14186--14186. Stable Stacks tag for an expressly allowed foundational or notational input; the counted construction remains the two concrete core obligations. The retained local core is not replaced by this pointer. See standard\_prerequisites.\par\medskip
\ifdefined\hypertarget\hypertarget{sourcepointer-12}{}\fi
\noindent\textit{Source reference.} Original section 0DPQ: focused-seeds/AGP/sources/stacks-snapshot-20261011/tags.txt, source lines 17160--17160. Stable Stacks tag for an expressly allowed foundational or notational input; the counted construction remains the two concrete core obligations. The retained local core is not replaced by this pointer. See standard\_prerequisites.\par\medskip
\ifdefined\hypertarget\hypertarget{sourcepointer-13}{}\fi
\noindent\textit{Source reference.} Original lemma 0DPR: focused-seeds/AGP/sources/stacks-snapshot-20261011/tags.txt, source lines 17161--17161. Stable Stacks tag for an expressly allowed foundational or notational input; the counted construction remains the two concrete core obligations. The retained local core is not replaced by this pointer. See standard\_prerequisites.\par\medskip
\ifdefined\hypertarget\hypertarget{sourcepointer-14}{}\fi
\noindent\textit{Source reference.} Original theorem 0D4X: focused-seeds/AGP/sources/stacks-snapshot-20261011/tags.txt, source lines 16467--16467. Stable Stacks tag for an expressly allowed foundational or notational input; the counted construction remains the two concrete core obligations. The retained local core is not replaced by this pointer. See standard\_prerequisites.\par\medskip
\ifdefined\hypertarget\hypertarget{sourcepointer-15}{}\fi
\noindent\textit{Source reference.} Original definition 01VH: focused-seeds/AGP/sources/stacks-snapshot-20261011/tags.txt, source lines 2446--2446. Stable Stacks tag for an expressly allowed foundational or notational input; the counted construction remains the two concrete core obligations. The retained local core is not replaced by this pointer. See standard\_prerequisites.\par\medskip
\ifdefined\hypertarget\hypertarget{sourcepointer-16}{}\fi
\noindent\textit{Source reference.} Original section 04SP: focused-seeds/AGP/sources/stacks-snapshot-20261011/tags.txt, source lines 6227--6227. Stable Stacks tag for an expressly allowed foundational or notational input; the counted construction remains the two concrete core obligations. The retained local core is not replaced by this pointer. See standard\_prerequisites.\par\medskip
\ifdefined\hypertarget\hypertarget{sourcepointer-17}{}\fi
\noindent\textit{Source reference.} Original section 03YL: focused-seeds/AGP/sources/stacks-snapshot-20261011/tags.txt, source lines 5188--5188. Stable Stacks tag for an expressly allowed foundational or notational input; the counted construction remains the two concrete core obligations. The retained local core is not replaced by this pointer. See standard\_prerequisites.\par\medskip
\ifdefined\hypertarget\hypertarget{sourcepointer-18}{}\fi
\noindent\textit{Source reference.} Original section 05X6: focused-seeds/AGP/sources/stacks-snapshot-20261011/tags.txt, source lines 7610--7610. Stable Stacks tag for an expressly allowed foundational or notational input; the counted construction remains the two concrete core obligations. The retained local core is not replaced by this pointer. See standard\_prerequisites.\par\medskip
\begin{lemma}
\label{lemma-polarized-diagonal-separated-fp}
The diagonal of $\Polarizedstack$ is separated
and of finite presentation.
\end{lemma}

\section{Properties of the stack of polarized proper schemes}
\label{section-polarized}

\noindent
In this section we discuss properties of the moduli stack
$$
\Polarizedstack \longrightarrow \Spec(\mathbf{Z})
$$
whose category of sections over a scheme $S$ is the category of
proper, flat, finitely presented scheme over $S$ endowed
with a relatively ample invertible sheaf. This is an algebraic
stack by Quot, Theorem \ref{quot-theorem-polarized-algebraic}.

\begin{situation}
\label{situation-polarized}
We define a category $\Polarizedstack$ as follows. Objects are
pairs $(X \to S, \mathcal{L})$ where
\begin{enumerate}
\item $X \to S$ is a morphism of schemes which is proper, flat, and
of finite presentation, and
\item $\mathcal{L}$ is an invertible $\mathcal{O}_X$-module
which is relatively ample on $X/S$
(Morphisms, Definition \ref{morphisms-definition-relatively-ample}).
\end{enumerate}
A morphism $(X' \to S', \mathcal{L}') \to (X \to S, \mathcal{L})$
between objects
is given by a triple $(f, g, \varphi)$ where $f : X' \to X$ and $g : S' \to S$
are morphisms of schemes which fit into a commutative diagram
$$
\xymatrix{
X' \ar[d] \ar[r]_f & X \ar[d] \\
S' \ar[r]^g & S
}
$$
inducing an isomorphism $X' \to S' \times_S X$, in other words, the
diagram is cartesian,
and $\varphi : f^*\mathcal{L} \to \mathcal{L}'$ is an isomorphism.
Composition is defined in the obvious manner (see
Examples of Stacks, Sections
\ref{examples-stacks-section-stack-of-spaces} and
\ref{examples-stacks-section-stack-of-quasi-coherent-sheaves}).
The forgetful functor
$$
p : \Polarizedstack \longrightarrow \Sch_{fppf},\quad
(X \to S, \mathcal{L}) \longmapsto S
$$
is how we view $\Polarizedstack$ as a category over $\Sch_{fppf}$
(see Section \ref{section-conventions} for notation).
\end{situation}

\begin{proof}
Recall that $\Polarizedstack$ is a limit preserving algebraic stack, see
Quot, Lemma \ref{quot-lemma-polarized-limits}.
By Limits of Stacks, Lemma \ref{stacks-limits-lemma-limit-preserving-diagonal}
this implies that
$\Delta : \Polarizedstack \to \Polarizedstack \times \Polarizedstack$
is limit preserving. Hence $\Delta$ is locally of finite presentation
by Limits of Stacks, Proposition
\ref{stacks-limits-proposition-characterize-locally-finite-presentation}.

\medskip\noindent
Let us prove that $\Delta$ is separated. To see this, it suffices to show
that given an affine scheme $U$ and two objects
$\upsilon = (Y, \mathcal{N})$ and $\chi = (X, \mathcal{L})$
of $\Polarizedstack$ over $U$, the algebraic
space
$$
\mathit{Isom}_{\Polarizedstack}(\upsilon, \chi)
$$
is separated. The rule which to an isomorphism $\upsilon_T \to \chi_T$
assigns the underlying isomorphism $Y_T \to X_T$ defines a morphism
$$
\mathit{Isom}_{\Polarizedstack}(\upsilon, \chi)
\longrightarrow
\mathit{Isom}_U(Y, X)
$$
Since we have seen in Lemmas \ref{lemma-Mor-s-lfp} and
\ref{lemma-Isom-in-Mor} that the target is
a separated algebraic space, it suffices to prove that this morphism
is separated. Given an isomorphism $f : Y_T \to X_T$
over some scheme $T/U$, then clearly
$$
\mathit{Isom}_{\Polarizedstack}(\upsilon, \chi)
\times_{\mathit{Isom}_U(Y, X), [f]} T
=
\mathit{Isom}(\mathcal{N}_T, f^*\mathcal{L}_T)
$$
Here $[f] : T \to \mathit{Isom}_U(Y, X)$ indicates the $T$-valued
point corresponding to $f$ and
$\mathit{Isom}(\mathcal{N}_T, f^*\mathcal{L}_T)$
is the algebraic space discussed in Section \ref{section-hom-isom}.
Since this algebraic space is affine over $U$, the claim implies
$\Delta$ is separated.

\medskip\noindent
To finish the proof we show that $\Delta$ is quasi-compact. Since
$\Delta$ is representable by algebraic spaces, it suffice to check
the base change of $\Delta$ by a surjective smooth morphism
$U \to \Polarizedstack \times \Polarizedstack$ is quasi-compact
(see for example Properties of Stacks, Lemma
\ref{stacks-properties-lemma-check-property-covering}).
We can assume $U = \coprod U_i$ is a disjoint union of affine opens.
Since $\Polarizedstack$ is limit preserving (see above), we
see that $\Polarizedstack \to \Spec(\mathbf{Z})$ is locally of
finite presentation, hence $U_i \to \Spec(\mathbf{Z})$ is
locally of finite presentation
(Limits of Stacks, Proposition
\ref{stacks-limits-proposition-characterize-locally-finite-presentation}
and Morphisms of Stacks, Lemmas
\ref{stacks-morphisms-lemma-composition-finite-presentation} and
\ref{stacks-morphisms-lemma-smooth-locally-finite-presentation}).
In particular, $U_i$ is Noetherian affine. This reduces us to the
case discussed in the next paragraph.

\medskip\noindent
In this paragraph, given a Noetherian affine scheme $U$ and two objects
$\upsilon = (Y, \mathcal{N})$ and $\chi = (X, \mathcal{L})$
of $\Polarizedstack$ over $U$, we show the algebraic space
$$
\mathit{Isom}_{\Polarizedstack}(\upsilon, \chi)
$$
is quasi-compact. Since the connected components of $U$ are open and closed
we may replace $U$ by these. Thus we may and do assume $U$ is connected.
Let $u \in U$ be a point. Let $P$ be the Hilbert polynomial
$n \mapsto \chi(Y_u, \mathcal{N}_u^{\otimes n})$, see
Varieties, Lemma \ref{varieties-lemma-numerical-polynomial-from-euler}.
Since $U$ is connected and since
the functions $u \mapsto \chi(Y_u, \mathcal{N}_u^{\otimes n})$
are locally constant (see 
Derived Categories of Schemes, Lemma
\ref{perfect-lemma-chi-locally-constant-geometric})
we see that we get the same Hilbert polynomial in every point of $U$.
Set
$\mathcal{M} = \text{pr}_1^*\mathcal{N}
\otimes_{\mathcal{O}_{Y \times_U X}} \text{pr}_2^*\mathcal{L}$
on $Y \times_U X$. Given
$(f, \varphi) \in \mathit{Isom}_{\Polarizedstack}(\upsilon, \chi)(T)$
for some scheme $T$ over $U$ then for every $t \in T$ we have
$$
\chi(Y_t, (\text{id} \times f)^*\mathcal{M}^{\otimes n}) =
\chi(Y_t,
\mathcal{N}_t^{\otimes n} \otimes_{\mathcal{O}_{Y_t}}
f_t^*\mathcal{L}_t^{\otimes n}) =
\chi(Y_t, \mathcal{N}_t^{\otimes 2n}) = P(2n)
$$
where in the middle equality we use the isomorphism
$\varphi : f^*\mathcal{L}_T \to \mathcal{N}_T$.
Setting $P'(t) = P(2t)$ we find that the morphism
$$
\mathit{Isom}_{\Polarizedstack}(\upsilon, \chi)
\longrightarrow
\mathit{Isom}_U(Y, X)
$$
(see earlier) has image contained in the intersection
$$
\mathit{Isom}_U(Y, X) \cap \mathit{Mor}^{P', \mathcal{M}}_U(Y, X)
$$
The intersection is an intersection of open subspaces of
$\mathit{Mor}_U(Y, X)$ (see Lemma \ref{lemma-Isom-in-Mor} and
Remark \ref{remark-Mor-numerical}).
Now $\mathit{Mor}^{P', \mathcal{M}}_U(Y, X)$
is a Noetherian algebraic space as it is of finite
presentation over $U$ by Lemma \ref{lemma-Mor-qc-over-base}.
Thus the intersection
is a Noetherian algebraic space too. Since the morphism
$$
\mathit{Isom}_{\Polarizedstack}(\upsilon, \chi)
\longrightarrow
\mathit{Isom}_U(Y, X) \cap \mathit{Mor}^{P', \mathcal{M}}_U(Y, X)
$$
is affine (see above) we conclude.
\end{proof}

\bibliographystyle{plain}
\bibliography{my}
\end{document}
